% A compact sequence view of the current application-owned dispatch workflow.
\begin{figure}[tbp]
\centering
\begin{tikzpicture}[
  actor/.style={draw,rounded corners=2pt,align=center,font=\small,
    minimum height=.7cm,minimum width=3.5cm,inner sep=5pt},
  message/.style={-{Stealth[length=1.8mm]},semithick},
  reply/.style={message,dashed},
  msg/.style={font=\footnotesize,align=center,fill=white,inner sep=2pt},
  note/.style={draw,rounded corners=2pt,font=\footnotesize,align=center,fill=white,inner sep=6pt}
]
\node[actor] (model) at (0,0) {Authoring model};
\node[actor] (app) at (5.3,0) {Application};
\node[actor] (checks) at (10.6,0) {MCP checks};
\foreach \x in {0,5.3,10.6} \draw[densely dashed,gray] (\x,-.4) -- (\x,-11.1);
\draw[message] (5.3,-1.2) -- node[msg,above]{Request, schema, catalogue} (0,-1.2);
\draw[reply] (0,-1.95) -- node[msg,above]{Proposal and check plan} (5.3,-1.95);
\node[note,text width=4.1cm] at (5.3,-2.85) {Validate candidate\\Freeze plan and arguments};
\draw[message] (5.3,-4.15) -- node[msg,above]{Run technical checks} (10.6,-4.15);
\draw[reply] (10.6,-4.9) -- node[msg,above]{Evidence or failure details} (5.3,-4.9);
\draw[message] (5.3,-5.65) -- node[msg,above]{Candidate and failure feedback} (0,-5.65);
\draw[reply] (0,-6.4) -- node[msg,above]{Complete replacement} (5.3,-6.4);
\node[note,text width=12.7cm] at (5.3,-7.35) {Recheck every replacement. Initial proposal plus at most three revisions.\\Skip selected difficulty checks until technical prerequisites pass.};
\draw[message] (5.3,-8.8) -- node[msg,above]{Assess difficulty if eligible} (10.6,-8.8);
\draw[reply] (10.6,-9.55) -- node[msg,above]{Difficulty result} (5.3,-9.55);
\node[note,text width=8.2cm] at (5.3,-10.5) {Save eligible artifact. Otherwise retain history\\Final status: checked, needs\_review or error};
\end{tikzpicture}
\caption{Authoring, checking, revision and finalisation sequence. Difficulty mismatch can also trigger revision within the same limit.}\label{fig:authoring}
\end{figure}
